\begin{figure}[tbp]
\centering
\begin{tikzpicture}[x=.94cm,y=.74cm,>=Stealth,font=\small]
 \foreach \i in {0,...,9}
   \node[font=\scriptsize] at (\i,.62) {\(\i\)};
 \node[anchor=east,font=\scriptsize] at (-.55,.62) {level};
 \foreach \sig/\col in {0/blue!65!black,1/green!45!black,2/orange!80!black}{
   \pgfmathsetmacro{\yy}{-1.18*\sig}
   \node[anchor=east] at (-.5,\yy) {\(\sigma=\sig\)};
   \draw[gray!50] (0,\yy)--(9,\yy);
   \foreach \i in {0,...,8}{
     \pgfmathtruncatemacro{\rr}{mod(\i,3)}
     \ifnum\rr=\sig\relax
       \draw[->,very thick,\col] (\i+.1,\yy)--(\i+.89,\yy);
     \else
       \draw[->,gray!45] (\i+.14,\yy)--(\i+.86,\yy);
     \fi
   }
   \foreach \i in {0,...,9}
     \fill[gray!65] (\i,\yy) circle (1.35pt);
   \pgfmathtruncatemacro{\first}{1+\sig}
   \draw[\col,thick] (0,\yy+.13)--(0,\yy+.32)
     --(\first,\yy+.32)--(\first,\yy+.13);
   \foreach \b in {0,1}{
     \pgfmathtruncatemacro{\aa}{\first+3*\b}
     \pgfmathtruncatemacro{\bb}{\aa+3}
     \draw[\col,thick] (\aa,\yy+.13)--(\aa,\yy+.32)
       --(\bb,\yy+.32)--(\bb,\yy+.13);
   }
 }
\end{tikzpicture}

\medskip
\begin{tikzpicture}[x=.94cm,y=1cm,font=\small,>=Stealth]
 \draw[thick] (0,0)--(9,0);
 \draw[thick] (0,-.12)--(0,.12) (9,-.12)--(9,.12);
 \node[above] at (4.5,.13) {block ancestor \(J\)};
 \fill[blue!12] (2.3,-1.35) rectangle (6,-.7);
 \draw[blue!65!black,thick] (2.3,-1.35) rectangle (6,-.7);
 \node at (3.85,-1.025) {immediate parent \(b\)};
 \fill[blue!65!black] (5.48,-1.35) rectangle (5.86,-.7);
 \node[below] at (5.67,-1.4) {\(v\)};
 \foreach \aa in {6.06,6.48}{
   \fill[orange!35] (\aa,-1.35) rectangle (\aa+.34,-.7);
   \draw[orange!80!black] (\aa,-1.35) rectangle (\aa+.34,-.7);
 }
 \node[anchor=west] at (7.15,-1.02) {neighbors \(v'\)};
 \draw[thin] (6.84,-1.02)--(7.08,-1.02);
 \draw[dashed] (5.48,-1.7)--(5.48,-2.4);
 \draw[dashed] (6.95,-1.5)--(6.95,-2.4);
 \draw[<->] (5.48,-2.27)--(6.95,-2.27);
 \node[anchor=west] at (7.05,-2.27) {\(\sim G^{-1/4}d(i)\)};
 \node[anchor=east] at (4.9,-2.27) {local neighbor radius};
\end{tikzpicture}
\caption{Shifted blocks and boundary masses, schematically.
The upper diagram uses \(K=3\). Each shift starts with \(1+\sigma\)
transitions and then uses blocks of \(K\) transitions; colored arrows
are the final transitions for that shift.
At a final transition \(i\to i+1\), the mass \(m_v\) includes all of
\(n_b\) and the masses of neighboring children outside \(b\) but inside
\(J\). The neighbor radius is determined by the immediate scale \(d(i)\),
even when \(J\) begins much earlier. Neither diagram is to scale.}
\label{fig:shifted-blocks}
\end{figure}
